\lecture{33}{Некапацитированное экономическое диспетчирование}

\sourcevideo
    {https://www.youtube.com/playlist?list=PLOzRYVm0a65fhKDS8cYIC7YCuVGJ4FOJx}
    {Week 9: Lecture 33: Algorithm for Uncapacitated EDP}

Рассматривается задача
\[
\begin{aligned}
    \min_{P_1,\ldots,P_N}\quad
    &\sum_{i=1}^{N}C_i(P_i),
    \\
    \text{при условии}\quad
    &\sum_{i=1}^{N}P_i=P_{\mathrm D},
\end{aligned}
\]
где
\[
    C_i(P_i)=\alpha_iP_i^2+\beta_iP_i+\gamma_i,
    \qquad \alpha_i>0.
\]
В оптимуме
\[
    2\alpha_iP_i^\star+\beta_i=\lambda^\star,
    \qquad
    \sum_{i=1}^{N}P_i^\star=P_{\mathrm D}.
\]
Следовательно, распределённый алгоритм должен сохранять баланс
мощности, установить локальную связь между мощностью и предельной
стоимостью и затем согласовать предельные стоимости.

\section{Локальные переменные и алгоритм}

Пусть генераторы соединены связным неориентированным графом
\[
    \mathcal G=(\mathcal V,\mathcal E),
    \qquad
    \mathcal V=\{1,\ldots,N\}.
\]
Через \(\mathcal N_i\) обозначим множество соседей генератора \(i\).
Для \(z\in\mathbb R\) и \(\mu>0\) положим
\[
    \operatorname{sig}^{\mu}(z)
    =
    |z|^\mu\operatorname{sign}(z).
\]
Эта функция нечётна и удовлетворяет равенству
\[
    z\operatorname{sig}^{\mu}(z)=|z|^{1+\mu}.
\]
Выберем
\[
    k_1,k_2,r_1,r_2>0,
    \qquad
    0<\mu_1<1<\mu_2.
\]

Каждый генератор хранит мощность \(P_i\) и локальную предельную
стоимость \(\lambda_i\). Ошибка локального условия оптимальности
равна
\[
    e_i
    =
    P_i-
    \frac{\lambda_i-\beta_i}{2\alpha_i}.
\]
Равенство \(e_i=0\) эквивалентно
\[
    \lambda_i=2\alpha_iP_i+\beta_i=C_i'(P_i).
\]
Определим консенсусный сигнал
\[
\begin{aligned}
    \Phi_i(\lambda)
    ={}&
    k_1
    \sum_{j\in\mathcal N_i}
    \operatorname{sig}^{\mu_1}(\lambda_j-\lambda_i)
    \\
    &+
    k_2
    \sum_{j\in\mathcal N_i}
    \operatorname{sig}^{\mu_2}(\lambda_j-\lambda_i).
\end{aligned}
\]
Распределённая динамика задаётся уравнениями
\[
    \dot P_i=\Phi_i(\lambda),
\]
\[
\begin{aligned}
    \dot\lambda_i
    =
    2\alpha_i
    \Bigl(
        \Phi_i(\lambda)
        &+r_1\operatorname{sig}^{\mu_1}(e_i)
        \\
        &+r_2\operatorname{sig}^{\mu_2}(e_i)
    \Bigr).
\end{aligned}
\]
Для реализации агенту нужны только собственные параметры
\(\alpha_i,\beta_i\), собственные переменные \(P_i,\lambda_i\) и
значения \(\lambda_j\) соседей.

\section{Сохранение баланса и динамика ошибки}

По нечётности знаковой степени и симметрии графа каждому слагаемому,
соответствующему ребру \(\{i,j\}\), отвечает противоположное
слагаемое. Поэтому
\[
    \sum_{i=1}^{N}\Phi_i(\lambda)=0.
\]
Следовательно,
\[
    \frac{\mathrm d}{\mathrm dt}
    \sum_{i=1}^{N}P_i=0
\]
и
\[
    \sum_{i=1}^{N}P_i(t)
    =
    \sum_{i=1}^{N}P_i(0).
\]
Если
\[
    \sum_{i=1}^{N}P_i(0)=P_{\mathrm D},
\]
то баланс мощности выполняется при всех \(t\ge0\).

Поскольку \(\beta_i\) постоянно,
\[
    \dot e_i
    =
    \dot P_i-
    \frac{\dot\lambda_i}{2\alpha_i}.
\]
После подстановки алгоритма консенсусные сигналы сокращаются:
\[
    \dot e_i
    =
    -r_1\operatorname{sig}^{\mu_1}(e_i)
    -r_2\operatorname{sig}^{\mu_2}(e_i).
\]
Таким образом, ошибки локальной оптимальности подчиняются автономной
динамике, не зависящей от графа и состояний соседей.

\section{Локальная оптимальность}

Рассмотрим функцию Ляпунова
\[
    V_e=\frac12\sum_{i=1}^{N}e_i^2.
\]
Её производная равна
\[
    \dot V_e
    =
    -r_1\sum_{i=1}^{N}|e_i|^{1+\mu_1}
    -r_2\sum_{i=1}^{N}|e_i|^{1+\mu_2}.
\]
Положим
\[
    p_1=\frac{1+\mu_1}{2},
    \qquad
    p_2=\frac{1+\mu_2}{2}.
\]
Тогда \(0<p_1<1<p_2\). Для неотрицательных чисел выполняются
оценки
\[
    \sum_{i=1}^{N}|e_i|^{1+\mu_1}
    \ge
    \left(
        \sum_{i=1}^{N}e_i^2
    \right)^{p_1},
\]
\[
    \sum_{i=1}^{N}|e_i|^{1+\mu_2}
    \ge
    N^{1-p_2}
    \left(
        \sum_{i=1}^{N}e_i^2
    \right)^{p_2}.
\]
Так как \(\sum_i e_i^2=2V_e\), получаем
\[
    \dot V_e
    \le
    -c_{e,1}V_e^{p_1}
    -c_{e,2}V_e^{p_2},
\]
где
\[
    c_{e,1}=r_1 2^{p_1},
    \qquad
    c_{e,2}=r_2N^{1-p_2}2^{p_2}.
\]
По критерию фиксированного времени все ошибки обращаются в нуль не
позднее момента
\[
    T_1
    \le
    \frac1{c_{e,1}(1-p_1)}
    +
    \frac1{c_{e,2}(p_2-1)}.
\]
Следовательно, при \(t\ge T_1\)
\[
    P_i
    =
    \frac{\lambda_i-\beta_i}{2\alpha_i},
\]
и это равенство сохраняется при дальнейшем движении.

\section{Консенсус предельных стоимостей}

После первого этапа
\[
    \frac{\dot\lambda_i}{2\alpha_i}
    =
    \dot P_i
    =
    \Phi_i(\lambda).
\]
Положим
\[
    w_i=\frac1{2\alpha_i}.
\]
Тогда
\[
    \frac{\mathrm d}{\mathrm dt}
    \sum_{i=1}^{N}w_i\lambda_i
    =
    \sum_{i=1}^{N}\Phi_i(\lambda)
    =0.
\]
В момент \(T_1\)
\[
\begin{aligned}
    \sum_{i=1}^{N}\frac{\lambda_i(T_1)}{2\alpha_i}
    &=
    \sum_{i=1}^{N}P_i(T_1)
    +
    \sum_{i=1}^{N}\frac{\beta_i}{2\alpha_i}
    \\
    &=
    P_{\mathrm D}
    +
    \sum_{i=1}^{N}\frac{\beta_i}{2\alpha_i}.
\end{aligned}
\]
Если \(\lambda_i\) сходятся к общему значению
\(\lambda_{\mathrm c}\), то
\[
    \lambda_{\mathrm c}
    \sum_{i=1}^{N}\frac1{2\alpha_i}
    =
    P_{\mathrm D}
    +
    \sum_{i=1}^{N}\frac{\beta_i}{2\alpha_i}.
\]
Отсюда
\[
    \lambda_{\mathrm c}
    =
    \frac{
        P_{\mathrm D}
        +
        \displaystyle\sum_{i=1}^{N}
        \frac{\beta_i}{2\alpha_i}
    }{
        \displaystyle\sum_{i=1}^{N}
        \frac1{2\alpha_i}
    }
    =
    \lambda^\star.
\]
Поэтому предельные мощности равны
\[
    P_i^\star
    =
    \frac{\lambda^\star-\beta_i}{2\alpha_i}.
\]
Консенсус достигается по предельным стоимостям, а мощности в общем
случае остаются различными.

\section{Фиксированное время второго этапа}

Положим
\[
    \widetilde\lambda_i=\lambda_i-\lambda^\star
\]
и введём диагональную матрицу
\[
    W
    =
    \operatorname{diag}
    \left(
        \frac1{2\alpha_1},\ldots,
        \frac1{2\alpha_N}
    \right).
\]
После момента \(T_1\)
\[
    \sum_{i=1}^{N}
    \frac{\widetilde\lambda_i}{2\alpha_i}=0.
\]
Рассмотрим
\[
    V_\lambda
    =
    \frac12
    \widetilde\lambda^{\mathsf T}
    W\widetilde\lambda.
\]
Для связного графа существует \(\kappa_W>0\), такое что на
указанном подпространстве
\[
    \widetilde\lambda^{\mathsf T}
    L\widetilde\lambda
    \ge
    \kappa_W
    \widetilde\lambda^{\mathsf T}
    W\widetilde\lambda.
\]
Используя ту же симметризацию по рёбрам и степенные неравенства, что
и в доказательстве фиксированно-временного консенсуса, получаем
\[
    \dot V_\lambda
    \le
    -c_{\lambda,1}V_\lambda^{p_1}
    -c_{\lambda,2}V_\lambda^{p_2}
\]
с некоторыми положительными константами
\(c_{\lambda,1}\) и \(c_{\lambda,2}\). Поэтому консенсус по
\(\lambda_i\) достигается за время
\[
    T_2
    \le
    \frac1{c_{\lambda,1}(1-p_1)}
    +
    \frac1{c_{\lambda,2}(p_2-1)},
\]
не зависящее от значений \(\lambda_i(T_1)\).

\begin{theorem}[Фиксированно-временное диспетчирование]
Пусть граф связен и неориентирован, функции стоимости квадратичны и
строго выпуклы, а начальные мощности удовлетворяют балансу
\[
    \sum_{i=1}^{N}P_i(0)=P_{\mathrm D}.
\]
Тогда приведённый распределённый алгоритм достигает единственного
решения некапацитированной задачи экономического диспетчирования не
позднее времени
\[
    T_{\max}\le T_1+T_2.
\]
При всех \(t\ge T_{\max}\)
\[
    P_i(t)=P_i^\star,
    \qquad
    \lambda_i(t)=\lambda^\star.
\]
\end{theorem}

\section{Интерпретация и ограничения}

Переменная \(\lambda_i\) является локальной предельной, или
инкрементальной, стоимостью. Классические дискретные методы,
согласовывающие такие переменные, называются алгоритмами согласования
инкрементальных стоимостей; рассмотренная схема является их
непрерывным фиксированно-временным вариантом.

Результат относится к некапацитированной модели. Ограничения
\[
    P_i^{\min}\le P_i\le P_i^{\max}
\]
алгоритмом автоматически не поддерживаются. Анализ также предполагает
квадратичные строго выпуклые стоимости, статический связный
неориентированный граф, симметричный и точный обмен, непрерывное
выполнение динамики, отсутствие потерь мощности и сетевых
ограничений. Прямая дискретизация Эйлера не обязана сохранять точную
фиксированно-временную гарантию. Капацитированная задача требует
изменения условий оптимальности и рассматривается далее.
